%% abstract.tex: abstract in PT and EN  (FEUP regulations)
%% -------------------------------------------------------

\chapter*{Resumo}
%
A oceanografia operacional depende de previsões e \emph{nowcasts} fiáveis para suportar a monitorização ambiental e a tomada de decisão. Contudo, observações \emph{in situ} continuam escassas devido ao custo e à cobertura limitada de campanhas com navios e de infraestruturas fixas. Veículos Autónomos Subaquáticos (AUVs) podem complementar redes de observação existentes ao fornecer medições persistentes e de alta resolução; ainda assim, limitações de autonomia energética, comunicações intermitentes e a variabilidade espaço--temporal dos processos oceânicos tornam a decisão de \emph{onde, quando e como} amostrar um problema sequencial sob incerteza. Em missões com múltiplos veículos, estes desafios ampliam-se pela necessidade de coordenação para evitar amostragem redundante, respeitando simultaneamente restrições operacionais de toda a frota.

Este trabalho enquadra o problema de amostragem \emph{dinâmica e multi-agente} em ambientes oceânicos tempo-variantes e incertos, com o objetivo de melhorar a estimação de campos e/ou a qualidade de previsão de variáveis como temperatura e salinidade. A hipótese central é que uma abordagem de \emph{amostragem consciente de regimes}---isto é, detetar regimes/estruturas distintas (por exemplo, massas de água, frentes, camadas de estratificação ou plumas fluviais) e condicionar decisões e atualizações do modelo a essa informação---permite recolher dados mais informativos e produzir modelos mais robustos do que estratégias que assumem homogeneidade. Em particular, espera-se que a consciência de regimes reduza a degradação de desempenho associada à mistura de observações recolhidas em regiões com dinâmicas muito diferentes.

A dissertação irá investigar um \emph{pipeline} em malha fechada no qual previsões e produtos de incerteza guiam o planeamento, as medições recolhidas alimentam uma etapa de atualização do modelo (ou um \emph{proxy} de atualização alinhado com o âmbito do trabalho) e o sistema replaneia iterativamente. O \emph{pipeline} integra: (i) inferência e agrupamento de regimes em tempo quase-real (online), recorrendo a modelos probabilísticos baseados em perfis (e.g., GMM/PCM), modelos sequenciais com comutação de regimes (e.g., HMM) e/ou representações compactas (embeddings ou \emph{sparse codes} obtidos por \emph{dictionary learning}); (ii) objetivos de planeamento informativo baseados em redução de incerteza e ganho de informação; (iii) coordenação multi-AUV sob restrições de comunicação, redundância e balanceamento de esforço; e (iv) custos/viabilidade conscientes de correntes em escoamentos oceânicos tempo-variantes.

Como resultado, espera-se produzir um enquadramento comparativo crítico da literatura relevante, uma proposta coerente de amostragem adaptativa multi-agente consciente de regimes, e um plano de validação com \emph{baselines}, métricas e cenários representativos para quantificar ganhos em informação, redução de redundância, custo operacional e desempenho de estimação/previsão.

\bigskip
\noindent\textbf{Palavras-chave:} amostragem adaptativa; veículos autónomos subaquáticos; coordenação multi-agente; deteção de regimes; planeamento informativo; correntes oceânicas tempo-variantes; atualização/assimilação de dados; aprendizagem automática.

\chapter*{Abstract}
%
Operational oceanography depends on reliable nowcasts and forecasts to support environmental monitoring and decision-making, yet in situ observations remain sparse due to the cost and limited coverage of ship-based campaigns and fixed observing infrastructure. Autonomous Underwater Vehicles (AUVs) can complement existing observing systems by providing persistent, high-resolution measurements; however, endurance constraints, intermittent communications, and the space--time variability of ocean processes make deciding \emph{where, when, and how} to sample a sequential decision problem under uncertainty. These challenges are amplified in multi-vehicle missions, where coordination is required to avoid redundant sampling while respecting fleet-wide operational constraints.

This work frames the problem of \emph{dynamic, multi-agent} sampling in time-varying and uncertain ocean environments, with the explicit goal of improving field estimation and/or forecast skill for variables such as temperature and salinity. The central hypothesis is that \emph{regime-aware sampling}---i.e., detecting distinct regimes/structures (e.g., water masses, fronts, stratification layers, or river plumes) and conditioning decisions and model updates on that information---enables more informative data collection and more robust modelling than approaches that assume environmental homogeneity. In particular, regime awareness is expected to mitigate performance degradation caused by mixing observations collected in regions governed by markedly different dynamics.

The dissertation will investigate a closed-loop pipeline in which forecasts and uncertainty products guide planning, collected measurements feed a model update stage (or a scope-aligned update proxy), and the system replans iteratively. The pipeline integrates: (i) near-real-time (online) regime inference and clustering using probabilistic profile-based models (e.g., GMM/PCM), sequential switching-regime models (e.g., HMMs), and/or learned compact representations (embeddings or sparse codes obtained via dictionary learning); (ii) informative planning objectives based on uncertainty reduction and information gain; (iii) multi-AUV coordination under communication, redundancy, and workload-balance constraints; and (iv) flow-aware feasibility and costs in time-varying ocean currents.

The expected outcome is a critical comparative framework of the relevant literature, a coherent regime-aware multi-agent adaptive sampling pipeline, and a validation plan with baselines, metrics, and representative scenarios to quantify improvements in information gain, redundancy reduction, operational cost, and forecast/estimation performance.

\bigskip
\noindent\textbf{Keywords:} adaptive sampling; autonomous underwater vehicles; multi-agent coordination; regime detection; informative planning; time-varying ocean flows; model update/data assimilation; machine learning.
